\documentclass[a4paper,10pt]{report}\input{../0.Préambule/Préambule_cours_prof}\begin{document}

\newpage
\section*{ALGORITHMIQUE ET PROGRAMMATION - Travaux pratiques 3}
\addcontentsline{toc}{section}{ALGORITHMIQUE ET PROGRAMMATION - Travaux pratiques 3}

\noindent \textbf{Objectif 1 :} \quad Frise Carrés

\medskip
\noindent \begin{minipage}{13.5cm}
\begin{dinglist}{111}
\item Crée un bloc nommé Carré qui permet de tracer un carré de côté $100$.
\item En utilisant les blocs, écris un programme permettant de tracer la frise ci-dessous.
\end{dinglist}
\begin{center}
\includegraphics[scale=0.5]{../40.Image/53.Frises_1.png}
\end{center}

\vspace{-5mm}
\ding{226} \textit{Appelle le professeur pour la vérification}
\end{minipage}
\begin{minipage}{5cm}
\begin{scratch}[scale=0.7]
\initmoreblocks{définir \namemoreblocks{Carré}}
\blockrepeat{répéter \ovalnum{4} fois}
{
\blockmove{avancer de \ovalnum{100}}
\blockmove{tourner \turnleft{} de \ovalnum{90} degrés}
}
\end{scratch}
\end{minipage}

\bigskip
\noindent \textbf{Objectif 2  :} \quad Frise Carrés espacés liés

\medskip
\noindent \begin{minipage}{9.5cm}
\begin{dinglist}{111}
\item Modifie le bloc Carré, puis écris un programme qui permet de tracer la frise ci-contre, composée de carrés de côtés $50$ et espacés de $20$ pixels chacun.
\end{dinglist}
\end{minipage}
\hspace{5mm}
\begin{minipage}{7cm}
\begin{center}
\includegraphics[scale=0.35]{../40.Image/53.Frises_2.png}
\end{center}
\end{minipage}

\bigskip
\noindent \textbf{Objectif 3  :} \quad Frise Carrés espacés

\medskip
\noindent \begin{minipage}{9.5cm}
\begin{dinglist}{111}
\item Écris un programme qui permet de tracer la frise ci-contre, composée de carrés de côtés $50$ et espacés de $20$ pixels chacun.
\end{dinglist}
\ding{226} \textit{Appelle le professeur pour la vérification}
\end{minipage}
\hspace{5mm}
\begin{minipage}{7cm}
\begin{center}
\includegraphics[scale=0.35]{../40.Image/53.Frises_3.png}
\end{center}
\end{minipage}

\bigskip
\noindent \textbf{Objectif 4 :} \quad Frise Carrés Rectangles

\medskip
\noindent \begin{minipage}{12.5cm}
\begin{dinglist}{111}
\item A la manière du bloc carré, crée un bloc nommé Rectangle qui permet de tracer un rectangle (vertical) de longueur $50$ et de largeur $20$.
\item En utilisant les blocs, écris un programme permettant de tracer la frise ci-contre.
\end{dinglist}
\end{minipage}
\hspace{5mm}
\begin{minipage}{4cm}
\begin{center}
\includegraphics[scale=0.4]{../40.Image/53.Frises_4.png}
\end{center}
\end{minipage}

\bigskip
\noindent \textbf{Objectif 5  :} \quad Frise Carrés Triangles

\medskip
\noindent \begin{minipage}{9.5cm}
\begin{dinglist}{111}
\item De la même manière, écris un programme permettant de tracer la frise ci-contre avec des triangles et des carrés de côté $50$.
\item \textit{As tu pensé à créer un bloc triangle ?}
\end{dinglist}
\end{minipage}
\hspace{5mm}
\begin{minipage}{7cm}
\begin{center}
\includegraphics[scale=0.45]{../40.Image/53.Frises_5.png}
\end{center}
\end{minipage}

\bigskip
\noindent \textbf{Objectif 6  :} \quad Frise Maison

\vspace{-2mm}
\noindent \begin{minipage}{9.5cm}
\begin{dinglist}{111}
\item Écris un programme permettant de tracer la frise ci-contre avec des triangles et des carrés de côté $50$.
\end{dinglist}
\ding{226} \textit{Appelle le professeur pour la vérification}
\end{minipage}
\hspace{5mm}
\begin{minipage}{7cm}
\begin{center}
\includegraphics[scale=0.4]{../40.Image/53.Frises_6.png}
\end{center}
\end{minipage}

\bigskip
\hrulefill
\bigskip

\noindent \textbf{Objectif bonus 1 :} \quad Frise Maison 2

\vspace{-4mm}
\noindent \begin{minipage}{10cm}
\vspace{3mm}
\begin{dinglist}{111}
\item Modifie le programme précédent pour obtenir la frise suivante. On pourra utiliser le bloc :
\begin{scratch}[scale=0.7]
\blockpen{ajouter \ovalnum{1} à la taille du stylo}
\end{scratch}
\end{dinglist}
\end{minipage}
\hspace{5mm}
\begin{minipage}{7cm}
\begin{center}
\includegraphics[scale=0.5]{../40.Image/53.Frises_7.png}
\end{center}
\end{minipage}

\bigskip
\noindent \begin{minipage}{7.5cm}
\textbf{Objectif bonus 2 :} \quad Frise Maison Emboitées

\medskip
\noindent \begin{minipage}{5cm}
\begin{dinglist}{111}
\item Écris un programme permettant de tracer cette construction. Tu peux créer la variable :
\begin{scratch}[scale=0.7]
\blockvariable{coté}
\end{scratch}
\end{dinglist}
\end{minipage}
\hspace{2mm}
\begin{minipage}{2cm}
\begin{center}
\includegraphics[scale=0.4]{../40.Image/53.Frises_8.png}
\end{center}
\end{minipage}
\end{minipage}
\begin{minipage}{8.5cm}
\textbf{Objectif bonus 3 :} \quad Frise Maison Emboitées

\medskip
\noindent \begin{minipage}{4cm}
\begin{dinglist}{111}
\item A partir de ce motif 
\includegraphics[scale=0.5]{../40.Image/53.Frises_9.png}, trace :

Tu peux créer le bloc :
\begin{scratch}[scale=0.7]
\blockmoreblocks{motif}
\end{scratch}
\end{dinglist}
\end{minipage}
\hspace{5mm}
\begin{minipage}{4cm}
\begin{center}
\includegraphics[scale=0.4]{../40.Image/53.Frises_10.png}
\end{center}
\end{minipage}
\end{minipage}
\end{document}